\documentclass[a4paper,12pt]{letter}

\usepackage{tikz}
\usepackage{amsmath}
\usepackage{amssymb}
\usepackage{nopageno}
\usepackage{amsmath}
%\usepackage{enumitem}
% \usepackage{mathtools}
%\usepackage{eqnarray,amsmath}
\usepackage[a4paper, total={7.5in, 11.25in}]{geometry}
%\usepackage[margin=2.5cm]{geometry}
\usepackage[utf8]{inputenc}
\usepackage[T1]{fontenc}
\usepackage{lmodern}
%\usepackage{amsmath, amssymb}
\usepackage{setspace}
\setstretch{1.2}
\usepackage{pgfplots}
\pgfplotsset{compat=1.18}

\usetikzlibrary{calc,3d,decorations.pathmorphing,decorations.pathreplacing,decorations.markings,snakes,arrows.meta,calc,patterns}
\tikzset{snake it/.style={decorate, decoration=snake}}

\begin{document}

\begin{center}
 (KTXFI1EBNF) Physics I.  Homework No.~4\\ \smallskip
May~23,~2026
 \end{center}

\begin{enumerate}

\item  Given the sum of two one-directional oscillation as the following $x(t) = 4\sin(\omega t) + 5\cos(\omega t).$ This describes the displacement - time function of the oscillation, where t represents time, $\omega$ denotes the circular frequency of the oscillation, and $\omega$ is a constant.
\begin{enumerate}
  \item[a)] Prove that the above x(t) oscillation, given by the displacement-time function, performs simple harmonic motion.
  \item[b)] What is the value of the resultant oscillation given by the above displacement-time function?
\end{enumerate}

% ----------------------------------------------------------------------------------------------------------------------------------------

\item Consider three one-directional oscillations with equal frequencies, whose amplitudes and initial phases are given as follows:
\begin{eqnarray*}
A_1 = 5 \cdot 10^{-3},\mathrm{m}&,& \varphi_1 = \pi/5 \\
A_2 = 2 \cdot 10^{-3}\,\mathrm{m}&,& \varphi_2 = \pi/4 \\
A_3 = 2 \cdot 10^{-3}\,\mathrm{m}&,& \varphi_3 = 13\pi/11
\end{eqnarray*}
Let’s combine these three oscillations. Calculate the amplitude and initial phase of the resultant vibration!

% ----------------------------------------------------------------------------------------------------------------------------------------

\item  Calculate and determine the curve described by the resultant of the perpendicular vibrations given by the equations:
\begin{eqnarray*}
x(t) &=& x_0 \cos(\omega t),\\  
y(t) &=& y_0 \sin(\omega t).  
\end{eqnarray*}


% ----------------------------------------------------------------------------------------------------------------------------------------

\item Let $f(\omega t \pm kx)$ be a function whose argument is at least twice differentiable, continuous, and arbitrary, where $\omega$: constant circular frequency, k : wavenumber (constant), $\varphi$ initial phase (constant), t : time, x spatial coordinate. Prove that the function f described above is a solution to the one-dimensional wave equation.

% ----------------------------------------------------------------------------------------------------------------------------------------

\item The displacement of a wave as a function of time is given by $$y(x,t) = 0.27 \sin(12x - 500t)\,mm,$$ where x is in meters and t is given in seconds. Determine at the point x = 0.2 meters in the case of a transverse wave:
\begin{enumerate}
\item[a)] its velocity,
\item[b)] its acceleration!
\item[c)] What is the displacement of the point at t = 4s seconds?
\end{enumerate}
  
% ----------------------------------------------------------------------------------------------------------------------------------------

\end{enumerate}

\bigskip

\rightline{Dr.~Gabor~FACSKO}
\rightline{\textit{facsko.gabor@uni-obuda.hu}}

\vfill

\end{document}
